\subsection{Measures and additive set functions of compact sets}

THEOREM 1. --- Let T be a topological space, and I a mapping of \(\mathfrak{K}(\mathrm{T})\) into \(\mathbf{R}_{+}\) . In order that there exist a measure \(\mu\) on T such that \(\mathrm{I}(\mathrm{K}) = \mu^{\bullet}(\mathrm{K})\) for all \(\mathrm{K} \in \mathfrak{K}(\mathrm{T})\) , it is necessary and sufficient that I satisfy the following conditions:

\begin{enumerate}
\def\labelenumi{\arabic{enumi})}
\item
  If K and L are compact subsets of T such that \(K \subset L\) , then \(I(K) \leqslant I(L)\) (`I is increasing').
\item
  If K and L are compact subsets of T, then I(K ∪ L) ≤ I(K) + I(L).
\item
  If K and L are disjoint compact subsets of T, then \(\mathrm{I}(\mathrm{K}\cup\mathrm{L})=\mathrm{I}(\mathrm{K})+\mathrm{I}(\mathrm{L})\) (`I is additive').
\item
  For every decreasing directed family \((\mathrm{K}_{\alpha})_{\alpha \in \mathrm{A}}\) of compact subsets of T, one has \(\mathrm{I}\big(\bigcap_{\alpha \in \mathrm{A}}\mathrm{K}_{\alpha}\big) = \inf_{\alpha \in \mathrm{A}}\mathrm{I}(\mathrm{K}_{\alpha})\) .
\item
  For every \(x \in \mathbf{T}\) , there exists a neighborhood \(\mathbf{V}\) of \(x\) such that
\end{enumerate}

\[
\sup_{\substack{\mathbf{K}\in \mathfrak{K}(\mathbf{T})\\ \mathbf{K}\subset \mathbf{V}}}\operatorname{I}(\mathbf{K}) <   + \infty
\]

(`I is locally bounded').

The measure \(\mu\) is then unique.

The uniqueness of \(\mu\) follows from the Cor. of Prop. 2 of §1, No.~2. The above conditions are necessary, the first three in obvious fashion, the last from the fact that \(\mu^{\bullet}\) is a locally bounded encumbrance, and condition 4) by the Cor. of Prop. 5 of §1, No.~6.

To show that these conditions are sufficient, we begin by treating the case that \(\mathbf{T}\) is compact.

Lemma 1. --- Assume that T is compact, and set \(l = \mathrm{I}(\mathrm{T})\) . For every \(\mathrm{A} \subset \mathrm{T}\) , set

\[
\mathrm{J}(\mathrm{A}) = \sup_{\substack{\mathrm{K}\in \mathfrak{K}(\mathrm{T})\\ \mathrm{K}\subset \mathrm{A}}}\mathrm{I}(\mathrm{K})\tag{1}
\]

and let \(\Phi\) be the set of \(A \subset T\) such that \(J(A) + J(\mathbb{C}A) = l\) . The set \(\Phi\) is then a clan that contains \(\mathfrak{K}(T)\) , and the function \(J\) on \(\Phi\) is increasing and additive.

It is clear that \(J\) is an increasing set function, extending \(I\) , and that \(J(A) + J(\mathbb{C}A) \leqslant l\) for every \(A \subset T\) .

Let K and S be two compact sets in T; we are first going to show that

\[
\mathrm{J} (\mathrm{K} \cap \mathrm{S}) + \mathrm{J} (\mathbb {C} \mathrm{K} \cap \mathrm{S}) = \mathrm{J} (\mathrm{S}).\tag{2}
\]

On considering the restrictions of I to \(\mathfrak{K}(\mathrm{S})\) and of J to \(\mathfrak{P}(\mathrm{S})\) , one reduces immediately to the case that \(\mathrm{S} = \mathrm{T}\) . Since T is normal, K is the intersection of the decreasing directed family of its compact neighborhoods, and condition 4) implies the existence, for every \(\varepsilon > 0\) , of a compact neighborhood H of K such that \(\mathrm{I}(\mathrm{H}) \leqslant \mathrm{I}(\mathrm{K}) + \varepsilon\) . Let L be the closure of T - H; L is compact, \(\mathrm{L} \cap \mathrm{K} = \varnothing\) and \(\mathrm{H} \cup \mathrm{L} = \mathrm{T}\) , therefore \(l = \mathrm{I}(\mathrm{H} \cup \mathrm{L}) \leqslant \mathrm{I}(\mathrm{H}) + \mathrm{I}(\mathrm{L}) \leqslant \mathrm{I}(\mathrm{K}) + \mathrm{I}(\mathrm{L}) + \varepsilon\) (condition 2)), whence the relation \(\mathrm{J}(\mathrm{K}) + \mathrm{J}(\mathbb{C}\mathrm{K}) \geqslant \mathrm{I}(\mathrm{K}) + \mathrm{I}(\mathrm{L}) \geqslant l - \varepsilon\) . Since \(\varepsilon\) is arbitrary, \(\mathrm{J}(\mathrm{K}) + \mathrm{J}(\mathbb{C}\mathrm{K}) = l\) . This proves the formula (2), as well as the inclusion \(\mathfrak{K}(\mathrm{T}) \subset \Phi\) .

Let us now prove that \(\Phi\) is a clan. Since \(\Phi\) is obviously stable under passage to the complement, it suffices to show that if \(A_{1}\) and \(A_{2}\) denote elements of \(\Phi\) , then \(A_{1} \cup A_{2} \in \Phi\) , or again that

\[
\mathrm{J} \left(\mathrm{A} _ {1} \cup \mathrm{A} _ {2}\right) + \mathrm{J} \left(\mathbf {C} \left(\mathrm{A} _ {1} \cup \mathrm{A} _ {2}\right)\right) \geqslant l.\tag{3}
\]

Denote by \(\varepsilon\) a number \(>0\) , and, for \(i = 1,2\) , let \(\mathbf{K}_i\) be a compact set contained in \(\mathbf{A}_i\) , and \(\mathbf{L}_i\) a compact set contained in \(\mathbf{C}\mathbf{A}_i\) , such that

\[
\mathrm{I} \left(\mathrm{K} _ {i}\right) \geqslant \mathrm{J} \left(\mathrm{A} _ {i}\right) - \varepsilon , \quad \mathrm{I} \left(\mathrm{L} _ {i}\right) \geqslant \mathrm{J} \left(\mathbb {C} \mathrm{A} _ {i}\right) - \varepsilon .
\]

Set \(\mathbf{M}_1 = \mathbf{K}_1\cup \mathbf{L}_1\) ; the relations \(l = J(\mathbf{M}_1) + \mathbf{J}(\mathbf{C}\mathbf{M}_1)\) and

\[
\mathrm{J} \left(\mathrm{M} _ {1}\right) = \mathrm{I} \left(\mathrm{K} _ {1}\right) + \mathrm{I} \left(\mathrm{L} _ {1}\right) \geqslant \mathrm{J} \left(\mathrm{A} _ {1}\right) + \mathrm{J} \left(\mathbf {C A} _ {1}\right) - 2 \varepsilon = l - 2 \varepsilon
\]

imply that \(\mathrm{J}(\mathbb{C}\mathrm{M}_1) \leqslant 2\varepsilon\) . Then if \(\mathrm{S}\) is a compact subset of \(\mathrm{T}\) , the relation (2) (applied to \(\mathrm{K} = \mathrm{M}_1\) ) implies \(\mathrm{J}(\mathrm{S}) \leqslant \mathrm{J}(\mathrm{M}_1 \cap \mathrm{S}) + 2\varepsilon\) , whence

\[
\mathrm{J} (\mathrm{S}) \leqslant \mathrm{J} (\mathrm{K} _ {1} \cap \mathrm{S}) + \mathrm{J} (\mathrm{L} _ {1} \cap \mathrm{S}) + 2 \varepsilon .
\]

Let us add the inequalities obtained by making \(S = K_{2}\) and \(S = L_{2}\) , and take into account the inequality \(\mathrm{J}(K_{2}) + \mathrm{J}(L_{2}) \geqslant l - 2\varepsilon\) and the fact that \(K_{1} \cap K_{2}\) , \(L_{1} \cap K_{2}\) and \(K_{1} \cap L_{2}\) are three disjoint compact sets contained in \(A_{1} \cup A_{2}\) . Denoting by C the union of these three compact sets, it follows that

\[
\begin{array}{r l} & {l - 2 \varepsilon \leqslant \mathrm{J} (\mathrm{K} _ {2}) + \mathrm{J} (\mathrm{L} _ {2}) \leqslant \mathrm{J} (\mathrm{C}) + \mathrm{J} (\mathrm{L} _ {1} \cap \mathrm{L} _ {2}) + 4 \varepsilon} \\ & {\qquad \leqslant \mathrm{J} (\mathrm{A} _ {1} \cup \mathrm{A} _ {2}) + \mathrm{J} \big (\mathbb {C} (\mathrm{A} _ {1} \cup \mathrm{A} _ {2}) \big) + 4 \varepsilon ,} \end{array}
\]

whence immediately the desired formula (3), in view of the arbitrariness of \(\varepsilon\) . This having been established, the preceding inequalities imply that \(\mathrm{J}(\mathbf{C}) \geqslant \mathrm{J}(\mathbf{A}_1 \cup \mathbf{A}_2) - 6\varepsilon\) ; if \(\mathbf{A}_1\) and \(\mathbf{A}_2\) are disjoint, then \(\mathbf{C}\) is the union of \(\mathbf{K}_1 \cap \mathbf{L}_2 \subset \mathbf{A}_1\) and \(\mathbf{K}_2 \cap \mathbf{L}_1 \subset \mathbf{A}_2\) , from which one deduces that \(\mathrm{J}(\mathbf{A}_1 \cup \mathbf{A}_2) \leqslant \mathrm{J}(\mathbf{A}_1) + \mathrm{J}(\mathbf{A}_2)\) . The reverse inequality being obvious, \(\mathbf{J}\) is indeed additive on \(\Phi\) , and the lemma is established.

Let us complete the proof of the theorem for the case that T is compact. Let \(\mathcal{E}(\Phi)\) be the vector space of \(\Phi\) -step functions on T, equipped with the topology of uniform convergence (Ch. IV, §4, No.~9, Def. 4); we shall again denote by J the positive linear form on \(\mathcal{E}(\Phi)\) associated with the additive function J (loc. cit., Prop. 18). Since \(J(T) = l\) , J is continuous and has norm \(l\) . Let then \(\mathcal{H}\) be the closure of \(\mathcal{E}(\Phi)\) for the topology of uniform convergence; one verifies at once that J may be extended by continuity to a positive linear form on \(\mathcal{H}\) , again denoted J. Since \(\mathcal{H}\) contains \(\mathcal{C}(T)\) (loc. cit., No.~10, Prop. 19) the restriction of J to \(\mathcal{C}(T)\) is a positive measure \(\mu\) . It remains to show that \(\mu^{\bullet}(K) = I(K)\) for every compact subset K of T. Now, we have \(\mu^{\bullet}(K) = \inf_{f \in S_K} \mu^{\bullet}(f)\) , where \(S_K\) denotes the set of elements of \(\mathcal{C}(T)\) that are \(\geqslant \varphi_K\) (§1, No.~6, Prop. 5). Since \(J(f) = \mu^{\bullet}(f)\) for \(f \in \mathcal{C}(T)\) , it clearly suffices to show that \(J(K) \geqslant \inf_{f \in S_K} J(f)\) . As in the proof of Lemma 1, let H be a compact neighborhood of K such that \(J(H) \leqslant J(K) + \varepsilon\) , and let \(f\) be a continuous function on T, between 0 and 1, equal to 1 on K and to 0 outside H (GT, IX, §4, No.~1, Prop. 1). Then

\[
\mathrm{J} (f) \leqslant \mathrm{J} (\mathrm{H}) \leqslant \mathrm{J} (\mathrm{K}) + \varepsilon ;
\]

\(\varepsilon\) being arbitrary, the desired inequality is proved, and the theorem is thus established when \(\mathbf{T}\) is compact.

Let us now pass to the general case. For every compact set L in T, let \(I_{L}\) be the restriction of I to \(\mathfrak{K}(L)\) . By the special case just treated, there exists a measure \(\mu_{L}\) on L, unique, such that \(\mu_{\mathrm{L}}(\mathrm{K}) = \mathrm{I}_{\mathrm{L}}(\mathrm{K})\) for every compact set \(K \subset L\) . Let then \(L'\) be a compact set contained in L; we have \((\mu_{\mathrm{L}})_{\mathrm{L}'}^{\bullet}(\mathrm{K}) = \mu_{\mathrm{L}}^{\bullet}(\mathrm{K}) = \mu_{\mathrm{L}'}^{\bullet}(\mathrm{K})\) for every compact set \(K \subset L'\) , therefore \(\mu_{\mathrm{L}'} = (\mu_{\mathrm{L}})_{\mathrm{L}'}\) ; the mapping \(\mu : L \mapsto \mu_{L}\) is a premeasure. The condition 5) expresses that \(\mu\) is a measure, and the relation \(\mathrm{I}(\mathrm{K}) = \mu^{\bullet}(\mathrm{K})\) for all compact \(\mathbf{K} \subset \mathbf{T}\) is obvious.

Remarks. --- 1) The condition 4) may be replaced, in the statement of Theorem 1, by the following condition (`right-continuity'):

4') For every \(\mathbf{K} \in \mathfrak{K}(\mathbf{T})\) and every \(\varepsilon > 0\) , there exists an open set \(\mathbf{U}\) containing \(\mathbf{K}\) , such that \(\mathrm{I}(\mathrm{H}) \leqslant \mathrm{I}(\mathrm{K}) + \varepsilon\) for every compact set \(\mathrm{H} \subset \mathbf{U}\) .

For, if \(\mu\) is a measure, the function \(I: K \mapsto \mu^{\bullet}(K)\) satisfies \(4'\) (§1, No.~9, Prop. 13). Conversely, suppose that I satisfies 1) and \(4'\) ; let us show that I then satisfies 4). With notations as in the statement of Theorem 1, choose an \(\varepsilon > 0\) and an open set U containing the compact set \(K = \bigcap_{\alpha \in A} K_{\alpha}\) and such that \(4'\) ) is

satisfied. There then exists an index \(\beta \in A\) such that \(K_{\beta} \subset U\) , and this implies

\[
\inf _ {\alpha \in \mathbf {A}} \mathrm{I} (\mathrm{K} _ {\alpha}) \leqslant \mathrm{I} (\mathrm{K} _ {\beta}) \leqslant \mathrm{I} (\mathrm{K}) + \varepsilon
\]

and 4) is indeed verified.

\begin{enumerate}
\def\labelenumi{\arabic{enumi})}
\setcounter{enumi}{1}
\tightlist
\item
  The set of conditions 2) and 3) may be replaced, in the statement of Theorem 1, by the following condition:
\end{enumerate}

If K and L are compact subsets of T, then

\[
\mathrm{I} (\mathrm{K} \cup \mathrm{L}) + \mathrm{I} (\mathrm{K} \cap \mathrm{L}) = \mathrm{I} (\mathrm{K}) + \mathrm{I} (\mathrm{L}).
\]

Indeed, this condition implies 2) and 3), and on the other hand

\[
\mu^ {\bullet} (\mathrm{K} \cup \mathrm{L}) = \mu^ {\bullet} (\mathrm{K} \cap \mathrm{L}) = \mu^ {\bullet} (\mathrm{K}) + \mu^ {\bullet} (\mathrm{L})
\]

for every measure \(\mu\) , by virtue of the relation \(\varphi_{\mathrm{K}\cup\mathrm{L}} + \varphi_{\mathrm{K}\cap\mathrm{L}} = \varphi_{\mathrm{K}} + \varphi_{\mathrm{L}}\) between the characteristic functions.

